% --- 04.Cartas.tex ---

\cleardoublepage
\chapter{Grimorio de Combate: Spells y Auras}

\begin{loreblock}
``No son hechizos — son protocolos. El Padre los grabó en nuestra sangre antes de cerrar los ojos.''
\end{loreblock}

\vspace{0.3cm}
\textit{Las cartas tienen formato TCG (88×63mm). Imprimirlas a doble faz permite tener el efecto completo al dorso. El cuadro de Costo ancla la esquina inferior izquierda. Recortar por el borde negro.}

\vspace{0.5cm}

% ============================================================
% GRILLA DE CARTAS — 2 columnas × 4 filas por página
% ============================================================
\begin{center}
\begin{tcbitemize}[
    raster columns=2,
    raster rows=4,
    raster column skip=0.5cm,
    raster row skip=0.4cm,
    enhanced,
    width=8.8cm,
    height=6.3cm,
    sharp corners,
    colframe=VoidBlack,
    colback=AbyssalBlue!8,
    boxed title style={colback=VoidPurple, coltext=white},
    fontupper=\footnotesize,
    halign upper=flush left,
    valign=top,
    raster width=18.5cm,
    raster force size,
    boxsep=2pt, left=2pt, right=2pt, top=0pt, bottom=2pt
]

% ==========================================
% CARTA 1: BARRERA ENTRÓPICA
% ==========================================
\tcbitem[title={\textbf{BARRERA ENTRÓPICA} \hfill \textcolor{white}{\small Reacción}}]

    {\fontsize{6.5}{7}\selectfont
    \textbf{Tipo:} Reacción \hfill \textbf{PE:} 1 \hfill \textbf{Unidad:} Cualquiera}
    \vspace{2pt}

    \begin{tcolorbox}[colback=ChromeGray!15, colframe=ChromeGray!40, height=2.5cm,
                      halign=center, valign=center, sharp corners, size=small, nobeforeafter]
        \textcolor{ChromeGray}{\fontsize{7}{8}\selectfont $\square$ [Ilustración: escudo de energía\\ absorbiendo un golpe enemigo]}
    \end{tcolorbox}

    \vfill

    \noindent
    \begin{minipage}[b]{0.22\linewidth}
        \begin{tcolorbox}[colback=BloodPurple, colframe=VoidBlack, sharp corners,
                          halign=center, valign=center, boxsep=0pt,
                          left=0pt, right=0pt, top=0pt, bottom=0pt,
                          nobeforeafter, height=1.9cm]
            \textcolor{white}{\fontsize{5.5}{6}\selectfont COSTO PE\\[3pt] \Large \textbf{1}}
        \end{tcolorbox}
    \end{minipage}%
    \hfill
    \begin{minipage}[b]{0.75\linewidth}
        \begin{tcolorbox}[colback=white, colframe=VoidPurple, sharp corners, boxrule=0.5pt,
                          left=2pt, right=2pt, top=2pt, bottom=2pt,
                          nobeforeafter, height=1.4cm]
            \fontsize{6}{7}\selectfont
            \textbf{Efecto:} Anula todo el daño de un ataque enemigo. La unidad sufre \textbf{1 Daño Directo} inevitable (ignora RD).
        \end{tcolorbox}
        \vspace{2pt}
        \begin{center}
            \fontsize{5.5}{6}\selectfont \itshape
            \textcolor{VoidPurple}{``Mejor sangrar por el Padre que por el enemigo.''}
        \end{center}
    \end{minipage}


% ==========================================
% CARTA 2: BALIZA DE ENJAMBRE
% ==========================================
\tcbitem[title={\textbf{BALIZA DE ENJAMBRE} \hfill \textcolor{white}{\small Aura}}]

    {\fontsize{6.5}{7}\selectfont
    \textbf{Tipo:} Aura — Mantenimiento (1 PE + 1 PV) \hfill \textbf{Unidad:} Mago}
    \vspace{2pt}

    \begin{tcolorbox}[colback=ChromeGray!15, colframe=ChromeGray!40, height=2.5cm,
                      halign=center, valign=center, sharp corners, size=small, nobeforeafter]
        \textcolor{ChromeGray}{\fontsize{7}{8}\selectfont $\square$ [Ilustración: balizas de luz\\ conectando unidades en red]}
    \end{tcolorbox}

    \vfill

    \noindent
    \begin{minipage}[b]{0.22\linewidth}
        \begin{tcolorbox}[colback=BloodPurple, colframe=VoidBlack, sharp corners,
                          halign=center, valign=center, boxsep=0pt,
                          left=0pt, right=0pt, top=0pt, bottom=0pt,
                          nobeforeafter, height=1.9cm]
            \textcolor{white}{\fontsize{5.5}{6}\selectfont COSTO PE\\[3pt] \Large \textbf{2}}
        \end{tcolorbox}
    \end{minipage}%
    \hfill
    \begin{minipage}[b]{0.75\linewidth}
        \begin{tcolorbox}[colback=white, colframe=VoidPurple, sharp corners, boxrule=0.5pt,
                          left=2pt, right=2pt, top=2pt, bottom=2pt,
                          nobeforeafter, height=1.4cm]
            \fontsize{6}{7}\selectfont
            \textbf{Aura 20cm:} Los ataques de aliados dentro del aura causan \textbf{Daño Explosivo} (+1d4). Mantenimiento al inicio de cada ronda: 1 PE + 1 PV.
        \end{tcolorbox}
        \vspace{2pt}
        \begin{center}
            \fontsize{5.5}{6}\selectfont \itshape
            \textcolor{VoidPurple}{``Extendiendo la red neuronal del Padre.''}
        \end{center}
    \end{minipage}


% ==========================================
% CARTA 3: MANTO OSCURO
% ==========================================
\tcbitem[title={\textbf{MANTO OSCURO} \hfill \textcolor{white}{\small Aura}}]

    {\fontsize{6.5}{7}\selectfont
    \textbf{Tipo:} Aura — Mantenimiento (1 FOCUS) \hfill \textbf{Unidad:} Héroes}
    \vspace{2pt}

    \begin{tcolorbox}[colback=ChromeGray!15, colframe=ChromeGray!40, height=2.5cm,
                      halign=center, valign=center, sharp corners, size=small, nobeforeafter]
        \textcolor{ChromeGray}{\fontsize{7}{8}\selectfont $\square$ [Ilustración: siluetas de unidades\\ acelerando entre sombras]}
    \end{tcolorbox}

    \vfill

    \noindent
    \begin{minipage}[b]{0.22\linewidth}
        \begin{tcolorbox}[colback=BloodPurple, colframe=VoidBlack, sharp corners,
                          halign=center, valign=center, boxsep=0pt,
                          left=0pt, right=0pt, top=0pt, bottom=0pt,
                          nobeforeafter, height=1.9cm]
            \textcolor{white}{\fontsize{5.5}{6}\selectfont COSTO PE\\[3pt] \Large \textbf{2}}
        \end{tcolorbox}
    \end{minipage}%
    \hfill
    \begin{minipage}[b]{0.75\linewidth}
        \begin{tcolorbox}[colback=white, colframe=VoidPurple, sharp corners, boxrule=0.5pt,
                          left=2pt, right=2pt, top=2pt, bottom=2pt,
                          nobeforeafter, height=1.4cm]
            \fontsize{6}{7}\selectfont
            \textbf{Aura 10cm:} Aliados ganan \textbf{+5cm de MOV por cada PA} gastado en moverse. Mantenimiento: 1 FOCUS al inicio de cada ronda.
        \end{tcolorbox}
        \vspace{2pt}
        \begin{center}
            \fontsize{5.5}{6}\selectfont \itshape
            \textcolor{VoidPurple}{``Viajamos por las sombras entre los átomos.''}
        \end{center}
    \end{minipage}


% ==========================================
% CARTA 4: PULSO NEURONAL
% ==========================================
\tcbitem[title={\textbf{PULSO NEURONAL} \hfill \textcolor{white}{\small Ofensivo}}]

    {\fontsize{6.5}{7}\selectfont
    \textbf{Tipo:} Hechizo Estándar (1 PA) \hfill \textbf{Unidad:} Mago}
    \vspace{2pt}

    \begin{tcolorbox}[colback=ChromeGray!15, colframe=ChromeGray!40, height=2.5cm,
                      halign=center, valign=center, sharp corners, size=small, nobeforeafter]
        \textcolor{ChromeGray}{\fontsize{7}{8}\selectfont $\square$ [Ilustración: onda de energía\\ psiónica expandiéndose desde el mago]}
    \end{tcolorbox}

    \vfill

    \noindent
    \begin{minipage}[b]{0.22\linewidth}
        \begin{tcolorbox}[colback=BloodPurple, colframe=VoidBlack, sharp corners,
                          halign=center, valign=center, boxsep=0pt,
                          left=0pt, right=0pt, top=0pt, bottom=0pt,
                          nobeforeafter, height=1.9cm]
            \textcolor{white}{\fontsize{5.5}{6}\selectfont COSTO PE\\[3pt] \Large \textbf{2}}
        \end{tcolorbox}
    \end{minipage}%
    \hfill
    \begin{minipage}[b]{0.75\linewidth}
        \begin{tcolorbox}[colback=white, colframe=VoidPurple, sharp corners, boxrule=0.5pt,
                          left=2pt, right=2pt, top=2pt, bottom=2pt,
                          nobeforeafter, height=1.4cm]
            \fontsize{6}{7}\selectfont
            \textbf{AoE 12cm radio:} Todas las unidades enemigas en el área deben tirar \textbf{Res. Mágica (1d6+RM) vs. 1d6+FOCUS}. Las que fallen quedan \textbf{Aturdidas} (pierden 1 PA en su próximo turno).
        \end{tcolorbox}
        \vspace{2pt}
        \begin{center}
            \fontsize{5.5}{6}\selectfont \itshape
            \textcolor{VoidPurple}{``El Padre susurra. Sus enemigos gritan.''}
        \end{center}
    \end{minipage}

\end{tcbitemize}
\end{center}

% ============================================================
% SEGUNDA PÁGINA DE CARTAS
% ============================================================
\newpage

\begin{center}
\begin{tcbitemize}[
    raster columns=2,
    raster rows=4,
    raster column skip=0.5cm,
    raster row skip=0.4cm,
    enhanced,
    width=8.8cm,
    height=6.3cm,
    sharp corners,
    colframe=VoidBlack,
    colback=AbyssalBlue!8,
    boxed title style={colback=VoidPurple, coltext=white},
    fontupper=\footnotesize,
    halign upper=flush left,
    valign=top,
    raster width=18.5cm,
    raster force size,
    boxsep=2pt, left=2pt, right=2pt, top=0pt, bottom=2pt
]

% ==========================================
% CARTA 5: VÓRTICE DE VACÍO
% ==========================================
\tcbitem[title={\textbf{VÓRTICE DE VACÍO} \hfill \textcolor{white}{\small Control}}]

    {\fontsize{6.5}{7}\selectfont
    \textbf{Tipo:} Ritual (TC — Turno Completo) \hfill \textbf{Unidad:} Mago}
    \vspace{2pt}

    \begin{tcolorbox}[colback=ChromeGray!15, colframe=ChromeGray!40, height=2.5cm,
                      halign=center, valign=center, sharp corners, size=small, nobeforeafter]
        \textcolor{ChromeGray}{\fontsize{7}{8}\selectfont $\square$ [Ilustración: vórtice de oscuridad\\ atrayendo unidades hacia el centro]}
    \end{tcolorbox}

    \vfill

    \noindent
    \begin{minipage}[b]{0.22\linewidth}
        \begin{tcolorbox}[colback=BloodPurple, colframe=VoidBlack, sharp corners,
                          halign=center, valign=center, boxsep=0pt,
                          left=0pt, right=0pt, top=0pt, bottom=0pt,
                          nobeforeafter, height=1.9cm]
            \textcolor{white}{\fontsize{5.5}{6}\selectfont COSTO PE\\[3pt] \Large \textbf{3}}
        \end{tcolorbox}
    \end{minipage}%
    \hfill
    \begin{minipage}[b]{0.75\linewidth}
        \begin{tcolorbox}[colback=white, colframe=VoidPurple, sharp corners, boxrule=0.5pt,
                          left=2pt, right=2pt, top=2pt, bottom=2pt,
                          nobeforeafter, height=1.4cm]
            \fontsize{6}{7}\selectfont
            \textbf{Ritual:} Crea un vórtice en un punto a 20cm. Durante 2 rondas, todas las unidades enemigas en 10cm del vórtice deben gastar \textbf{+1 PA extra} para alejarse de él.
        \end{tcolorbox}
        \vspace{2pt}
        \begin{center}
            \fontsize{5.5}{6}\selectfont \itshape
            \textcolor{VoidPurple}{``El vacío no atrae. Simplemente consume.''}
        \end{center}
    \end{minipage}


% ==========================================
% CARTA 6: SEÑAL DE RASTREO
% ==========================================
\tcbitem[title={\textbf{SEÑAL DE RASTREO} \hfill \textcolor{white}{\small Soporte}}]

    {\fontsize{6.5}{7}\selectfont
    \textbf{Tipo:} Hechizo Estándar (1 PA) \hfill \textbf{Unidad:} Cualquiera}
    \vspace{2pt}

    \begin{tcolorbox}[colback=ChromeGray!15, colframe=ChromeGray!40, height=2.5cm,
                      halign=center, valign=center, sharp corners, size=small, nobeforeafter]
        \textcolor{ChromeGray}{\fontsize{7}{8}\selectfont $\square$ [Ilustración: marcador de objetivo\\ brillando sobre una unidad enemiga]}
    \end{tcolorbox}

    \vfill

    \noindent
    \begin{minipage}[b]{0.22\linewidth}
        \begin{tcolorbox}[colback=BloodPurple, colframe=VoidBlack, sharp corners,
                          halign=center, valign=center, boxsep=0pt,
                          left=0pt, right=0pt, top=0pt, bottom=0pt,
                          nobeforeafter, height=1.9cm]
            \textcolor{white}{\fontsize{5.5}{6}\selectfont COSTO PE\\[3pt] \Large \textbf{1}}
        \end{tcolorbox}
    \end{minipage}%
    \hfill
    \begin{minipage}[b]{0.75\linewidth}
        \begin{tcolorbox}[colback=white, colframe=VoidPurple, sharp corners, boxrule=0.5pt,
                          left=2pt, right=2pt, top=2pt, bottom=2pt,
                          nobeforeafter, height=1.4cm]
            \fontsize{6}{7}\selectfont
            \textbf{Rango 25cm:} Marca una unidad enemiga. Todos los aliados Biskrorum ganan \textbf{+1 a tiradas de Éxito} contra esa unidad hasta el final de la ronda. 1× por ronda.
        \end{tcolorbox}
        \vspace{2pt}
        \begin{center}
            \fontsize{5.5}{6}\selectfont \itshape
            \textcolor{VoidPurple}{``Objetivos compartidos. Voluntad compartida.''}
        \end{center}
    \end{minipage}


% ==========================================
% CARTA 7: PROTOCOLO DE AUTORREPARACIÓN
% ==========================================
\tcbitem[title={\textbf{AUTORREPARACIÓN} \hfill \textcolor{white}{\small Soporte}}]

    {\fontsize{6.5}{7}\selectfont
    \textbf{Tipo:} Hechizo Estándar (1 PA) \hfill \textbf{Unidad:} Cualquiera con exoesqueleto}
    \vspace{2pt}

    \begin{tcolorbox}[colback=ChromeGray!15, colframe=ChromeGray!40, height=2.5cm,
                      halign=center, valign=center, sharp corners, size=small, nobeforeafter]
        \textcolor{ChromeGray}{\fontsize{7}{8}\selectfont $\square$ [Ilustración: nanobots reparando\\ el metal fracturado de un exoesqueleto]}
    \end{tcolorbox}

    \vfill

    \noindent
    \begin{minipage}[b]{0.22\linewidth}
        \begin{tcolorbox}[colback=BloodPurple, colframe=VoidBlack, sharp corners,
                          halign=center, valign=center, boxsep=0pt,
                          left=0pt, right=0pt, top=0pt, bottom=0pt,
                          nobeforeafter, height=1.9cm]
            \textcolor{white}{\fontsize{5.5}{6}\selectfont COSTO PE\\[3pt] \Large \textbf{2}}
        \end{tcolorbox}
    \end{minipage}%
    \hfill
    \begin{minipage}[b]{0.75\linewidth}
        \begin{tcolorbox}[colback=white, colframe=VoidPurple, sharp corners, boxrule=0.5pt,
                          left=2pt, right=2pt, top=2pt, bottom=2pt,
                          nobeforeafter, height=1.4cm]
            \fontsize{6}{7}\selectfont
            \textbf{Auto:} Recupera \textbf{1d6 PV} en la zona con menos vida. Si una extremidad estaba incapacitada y recupera al menos 1 PV, se reactiva temporalmente por 1 ronda.
        \end{tcolorbox}
        \vspace{2pt}
        \begin{center}
            \fontsize{5.5}{6}\selectfont \itshape
            \textcolor{VoidPurple}{``No curamos. Reconfiguramos.''}
        \end{center}
    \end{minipage}


% ==========================================
% CARTA 8: COLAPSO DE BIOMASA
% ==========================================
\tcbitem[title={\textbf{COLAPSO DE BIOMASA} \hfill \textcolor{white}{\small Ofensivo}}]

    {\fontsize{6.5}{7}\selectfont
    \textbf{Tipo:} Reacción \hfill \textbf{Unidad:} Cualquiera — se activa al morir}
    \vspace{2pt}

    \begin{tcolorbox}[colback=ChromeGray!15, colframe=ChromeGray!40, height=2.5cm,
                      halign=center, valign=center, sharp corners, size=small, nobeforeafter]
        \textcolor{ChromeGray}{\fontsize{7}{8}\selectfont $\square$ [Ilustración: explosión de ácido\\ y biomasa al morir una unidad]}
    \end{tcolorbox}

    \vfill

    \noindent
    \begin{minipage}[b]{0.22\linewidth}
        \begin{tcolorbox}[colback=BloodPurple, colframe=VoidBlack, sharp corners,
                          halign=center, valign=center, boxsep=0pt,
                          left=0pt, right=0pt, top=0pt, bottom=0pt,
                          nobeforeafter, height=1.9cm]
            \textcolor{white}{\fontsize{5.5}{6}\selectfont COSTO PE\\[3pt] \Large \textbf{0}}
        \end{tcolorbox}
    \end{minipage}%
    \hfill
    \begin{minipage}[b]{0.75\linewidth}
        \begin{tcolorbox}[colback=white, colframe=VoidPurple, sharp corners, boxrule=0.5pt,
                          left=2pt, right=2pt, top=2pt, bottom=2pt,
                          nobeforeafter, height=1.4cm]
            \fontsize{6}{7}\selectfont
            \textbf{Al morir:} Esta unidad explota en un radio de \textbf{6cm}. Todas las unidades en el radio reciben \textbf{1d6 daño} (ignora RD). Tirar localización normal para cada una. No afecta a unidades Biskrorum.
        \end{tcolorbox}
        \vspace{2pt}
        \begin{center}
            \fontsize{5.5}{6}\selectfont \itshape
            \textcolor{VoidPurple}{``La muerte es el último disparo.''}
        \end{center}
    \end{minipage}

\end{tcbitemize}
\end{center}

% ============================================================
% TABLA RESUMEN DE CARTAS
% ============================================================
\newpage
\section*{Resumen de Cartas — Referencia Rápida}

\begin{tabularx}{\linewidth}{|l|c|c|c|X|}
\hline
\rowcolor{VoidPurple!25} \textbf{Nombre} & \textbf{Tipo} & \textbf{PE} & \textbf{PA} & \textbf{Efecto Breve} \\ \hline
Barrera Entrópica & Reacción & 1 & 2 & Anula 1 ataque, sufre 1 daño directo \\ \hline
Baliza de Enjambre & Aura & 2 & 1 & Aura 20cm: +1d4 daño explosivo aliados \\ \hline
Manto Oscuro & Aura & 2 & 1 & Aura 10cm: +5cm MOV por PA en movimiento \\ \hline
Pulso Neuronal & Ofensivo & 2 & 1 & AoE 12cm: Aturdido a los que fallen RM \\ \hline
Vórtice de Vacío & Control & 3 & TC & +1 PA para alejarse del vórtice por 2 rondas \\ \hline
Señal de Rastreo & Soporte & 1 & 1 & Marca enemigo: +1 ataque aliados vs. él \\ \hline
Autorreparación & Soporte & 2 & 1 & Recupera 1d6 PV en zona más dañada \\ \hline
Colapso de Biomasa & Reacción & 0 & — & Al morir: 1d6 daño en 6cm radio \\ \hline
\end{tabularx}

\vspace{0.5cm}

\begin{reglabiskrida}{Regla: Límite de Spells por Unidad}
\begin{itemize}
    \item Cada unidad puede equipar entre \textbf{1 y 6 cartas de spell} según el número de ranuras de Spells indicado en su ficha.
    \item Las cartas deben elegirse al armar la lista, antes de la partida.
    \item Los costos de PE de las cartas se pagan de las fuentes normales: reserva propia, Gema, o extracción del entorno (creando Zona Seca).
    \item Las cartas con etiqueta \textbf{Reacción} no cuestan PA adicional si están en la lista, pero sí cuestan PE.
    \item Las cartas con etiqueta \textbf{Ritual (TC)} consumen \textbf{toda la activación} de la unidad — no puede realizar ninguna otra acción ese turno.
\end{itemize}
\end{reglabiskrida}

\vspace{0.4cm}

\begin{tcolorbox}[colback=GoldRule!10, colframe=GoldRule,
    title={\textbf{CARTAS PENDIENTES DE DISEÑO}}]
\small Las siguientes cartas están en proceso de diseño y se agregarán en la próxima versión:
\begin{itemize}
    \item \textit{Invocación de Autómata} — Convoca un Nódulo adicional al tablero (Ritual, 4 PE)
    \item \textit{Enlace de Sacrificio} — Transfiere PV de una unidad a otra aliada (Soporte, 1 PE)
    \item \textit{Protocolo Berserker} — Elimina el límite de PA por 1 turno, luego la unidad queda Agotada permanentemente hasta el inicio de la siguiente ronda (Ofensivo, 3 PE)
    \item \textit{Escudo Colmena} — Invoca a 2 Nódulos frente a una unidad aliada como escudos vivientes (Reacción, 2 PE)
\end{itemize}
\end{tcolorbox}
